% ------------------------------- Chapter 2 -------------------------------
\chapter{1-D Mean-Line Design}

The 1-D mean-line design is the simplified modelling strategy used in the early
stages of turbine design. The flow is assumed axisymmetric and is averaged along
the radial direction, so that the whole stage is represented by the conditions on
the mean streamline. With this simplification the stage power, efficiency and
blade loading can be evaluated quickly and the design coefficients iterated.

The stage is divided into three stations: station~1 (stator inlet), station~2
(stator exit / rotor inlet) and station~3 (rotor exit).

% =====================================================================
\section{Input Data}

The following quantities are imposed by the project statement:

\begin{itemize}
    \item Working fluid: calorically perfect gas, $\gamma = 4/3 \approx 1.33$,
          $R = \SI{287}{J/kg/K}$.
    \item Turbine mass flow $\dot m = \SI{750}{kg/s}$; fuel-to-air ratio
          $f_{ar}=0.025$.
    \item Stage (extracted) shaft power $P = \SI{330}{MW}$ (the turbine also
          drives the compressor).
    \item Overall compressor pressure ratio $\Pi_C = 20$.
    \item Compressor admission $T_{adm} = \SI{15}{\degreeCelsius}$,
          $P_{adm} = \SI{101325}{Pa}$, isentropic efficiency $\eta_{is,C}=0.92$.
    \item Combustion-chamber Mach number $M_{comb}=0.15$.
    \item Turbine inlet static temperature (F-class) $TIT = \SI{1400}{\degreeCelsius}$.
    \item Inlet yaw angle $\alpha_1 = 0\degr$ (fully axial); rotational speed
          $N_{rot}=\SI{3000}{rpm}$ (synchronous generator).
    \item Blade peripheral speed limited to $U_m = \omega R_m \le \SI{350}{m/s}$;
          mechanical constraint $AN^2 < \SI{3e7}{m^2\,rpm^2}$.
    \item Constant mean radius $R_{mean}$ across the stage; fixed blading
          efficiencies $\eta_{S,R}=0.90$.
\end{itemize}

% =====================================================================
\section{Design Procedure}

\subsection{Inlet Condition Calculations}

The specific heats are obtained from $\gamma$ and $R$:
\begin{equation}
    c_p = \frac{\gamma R}{\gamma-1} = \SI{1156.70}{J/kg/K},\qquad
    c_v = \frac{R}{\gamma-1} = \SI{869.70}{J/kg/K}.
\end{equation}

Only air flows through the compressor, so its mass flow differs from the turbine:
\begin{equation}
    \dot m_c = \frac{\dot m_T}{1.025} = \SI{731.71}{kg/s}.
\end{equation}
With $\gamma_C = 1.4$ for the compressor, the absorbed compressor power is
\begin{equation}
    P_c = \frac{1}{\eta_{is}\,\eta_m}\,\dot m_c\,c_{p,C}\,T_{adm}
          \left[\Pi_C^{\frac{\gamma_C-1}{\gamma_C}}-1\right] = \SI{327.82}{MW},
\end{equation}
taking a mechanical efficiency $\eta_m = 0.95$. The compressor exit (turbine
inlet) static pressure follows from the pressure ratio,
\begin{equation}
    P_1 = \Pi_C \, P_{adm} = \SI{2026500}{Pa}.
\end{equation}
The turbine inlet stagnation properties are obtained from the combustion-chamber
Mach number:
\begin{align}
    P_{01} &= P_1\left(1+\tfrac{\gamma-1}{2}M_{comb}^2\right)^{\frac{\gamma}{\gamma-1}}
            = \SI{2056992}{Pa},\\
    T_{01} &= T_1\left(1+\tfrac{\gamma-1}{2}M_{comb}^2\right)
            = \SI{1679.21}{K},
\end{align}
with $T_1 = TIT = \SI{1673}{K}$.

\subsection{Required \texorpdfstring{$P_3$}{P3} Calculations}

The total power to be produced by the four-stage turbine is the sum of the net
shaft power and the compressor power, so the power per stage is
\begin{equation}
    P_{stage} = \frac{P + P_c}{4} = \SI{164.46}{MW}.
\end{equation}
The corresponding real specific work and isentropic enthalpy drop
(with $\eta_{TT}=0.87$) are
\begin{equation}
    \Delta H_{real} = \frac{P_{stage}}{\dot m} = \SI{219.27}{kJ/kg},\qquad
    \Delta H_{is} = \frac{\Delta H_{real}}{\eta_{TT}} = \SI{243.64}{kJ/kg}.
\end{equation}
The isentropic exit stagnation state and stage pressure ratio are
\begin{equation}
    T_{03,is} = T_{01}-\frac{\Delta H_{is}}{c_p} = \SI{1468.58}{K},\quad
    P_{03,is} = \SI{1198493}{Pa},\quad
    \Pi_{TT,is} = \frac{P_{01}}{P_{03,is}} \approx 1.72 .
\end{equation}
Assuming an exit Mach number $M_3 = 0.3$ (within the recommended $0.3$--$0.4$
band) closes the static exit state and gives a first estimate of $P_3$, which is
later refined by the iteration of Section~\ref{sec:enth}.

\subsection{Assumptions of Design Parameters}

Three coefficients are chosen and iterated to meet the design guidelines:

\paragraph{Stage loading $\psi = \Delta H_{real}/U^2$.}
For a high-pressure stage $\psi$ is typically $1.5$--$2$. It controls the width
of the velocity triangles, the turning and the exit Mach number. Fixing
$\psi = 1.8$ gives a blade speed
\begin{equation}
    U_m = \sqrt{\frac{\Delta H_{real}}{\psi}} = \SI{349.03}{m/s},
\end{equation}
just within the \SI{350}{m/s} mechanical limit.

\paragraph{Flow coefficient $\phi = V_x/U$.}
The flow coefficient sets the height of the velocity triangles. Using a Smith
chart~\cite{smith} at the chosen efficiency and $\psi$, a value $\phi = 0.5$ is
selected, giving an axial velocity $V_x = \SI{174.5}{m/s}$.

\paragraph{Degree of reaction $r = (h_2-h_3)/(h_1-h_3)$.}
For a cooled high-pressure stage the reaction is kept low, $0.35$--$0.5$; here
$r = 0.35$ is adopted.

% =====================================================================
\section{Thermodynamic Quantities}
\label{sec:thermo}

With the inlet state and the three coefficients fixed, the thermodynamic and
kinematic quantities are computed at every station. Station~2 follows from the
stator expansion ($V_{x2}=U\phi$, $P_2 = r(P_1-P_3)+P_3$, stator efficiency
$\eta_S$); the relative frame ($W_2$, $\beta_2$, $T_{02R}$, $P_{02R}$) is then
formed, and the rotor exit (station~3) is obtained from the relative isentropic
expansion to $P_3$ with rotor efficiency $\eta_R$. Table~\ref{tab:stations}
collects the converged values.

\begin{table}[H]
\centering
\caption{Thermodynamic and kinematic quantities at the three stations.}
\label{tab:stations}
\small
\begin{tabular}{l r r r}
\toprule
\textbf{Quantity} & \textbf{Station 1} & \textbf{Station 2} & \textbf{Station 3}\\
\midrule
Static pressure $P$ [Pa]            & 2026500 & 1428997 & 1107265 \\
Total pressure (abs.) $P_0$ [Pa]    & 2056992 & 1954202 & 1166813 \\
Rel. total pressure $P_{0R}$ [Pa]   & --      & 1521338 & 1472070 \\
Static temperature $T$ [K]          & 1673.0  & 1553.7  & 1470.4 \\
Total temperature (abs.) $T_0$ [K]  & 1679.2  & 1679.2  & 1489.6 \\
Rel. total temperature $T_{0R}$ [K] & --      & 1578.1  & 1578.1 \\
Absolute velocity $V$ [m/s]         & 174.5   & 538.8   & 211.0 \\
Axial velocity $V_x$ [m/s]          & 174.5   & 174.5   & 174.5 \\
Tangential velocity $V_u$ [m/s]     & 0.0     & 509.7   & 118.5 \\
Relative velocity $W$ [m/s]         & --      & 237.2   & 499.1 \\
Yaw angle $\alpha$ [\degr]          & 0.00    & 71.10   & 34.18 \\
Relative angle $\beta$ [\degr]      & --      & 42.64   & 69.53 \\
Absolute Mach $M$ [-]               & 0.22    & 0.70    & 0.30 \\
Relative Mach $M_w$ [-]             & --      & 0.31    & 0.67 \\
\bottomrule
\end{tabular}
\end{table}

The stator-exit Mach $M_2 = 0.70$ and the relative exit Mach $M_{w3}=0.67$ are
consistent with a low-reaction high-pressure stage in which the bulk of the
acceleration occurs in the stator.

% =====================================================================
\section{Velocity Triangles}

The velocity triangles at stations~2 and~3 are built from the blade speed $U$,
the absolute velocity $V$ and the relative velocity $W$, together with the flow
angle $\alpha$ and the blade angle $\beta$. They are plotted in
Fig.~\ref{fig:tri} and summarised in Table~\ref{tab:tri}.

\begin{figure}[H]
\centering
\includegraphics[width=0.8\textwidth]{Velocity Triangles at Rotor Inlet (2) and Exit (3).png}
\caption{Velocity triangles at rotor inlet (station 2) and exit (station 3).}
\label{fig:tri}
\end{figure}

\begin{table}[H]
\centering
\caption{Velocity-triangle parameters.}
\label{tab:tri}
\begin{tabular}{l r r}
\toprule
\textbf{Parameter} & \textbf{Station 2} & \textbf{Station 3}\\
\midrule
$U$ -- blade speed [m/s]        & 349.03 & 349.03 \\
$V$ -- absolute velocity [m/s]  & 538.78 & 210.95 \\
$W$ -- relative velocity [m/s]  & 237.24 & 499.05 \\
$\alpha$ -- flow angle [\degr]  & 71.10  & 34.18 \\
$\beta$ -- blade angle [\degr]  & 42.64  & 69.53 \\
\bottomrule
\end{tabular}
\end{table}

The total turning of the rotor is $\Delta\beta = \beta_2+\beta_3 = 112.2\degr$,
typical of a heavily loaded high-pressure rotor.

% =====================================================================
\section{h--s Diagram}

The enthalpy--entropy diagram (Fig.~\ref{fig:hs}) summarises the expansion: the
isentropic, static and stagnation enthalpies at each station and the kinetic
energy terms $V^2/2$ and $W^2/2$. The enthalpy ladder is listed in
Table~\ref{tab:hs}.

\begin{table}[H]
\centering
\caption{Enthalpy ladder of the stage (J/kg).}
\label{tab:hs}
\begin{tabular}{l r l r}
\toprule
\textbf{Quantity} & \textbf{Value} & \textbf{Quantity} & \textbf{Value}\\
\midrule
$H_{01}=H_{02}$ & 1942338 & $h_3$    & 1700813 \\
$h_1$           & 1935154 & $h_{3s}$ & 1686977 \\
$h_2$           & 1797198 & $h_{3ss}$& 1665655 \\
$h_{2s}$        & 1781071 & $H_{03}$ & 1723063 \\
$H_{0R}$        & 1825338 & $H_{03ss}$& 1687906 \\
\bottomrule
\end{tabular}
\end{table}

\begin{figure}[H]
\centering
\includegraphics[width=0.9\textwidth]{h-s Diagram for Turbine Stage.png}
\caption{h--s diagram of the first turbine stage.}
\label{fig:hs}
\end{figure}

\subsection{Comparison of Enthalpy (closure of the design)}
\label{sec:enth}

The exit static pressure $P_3$ is not known a priori. The stage enthalpy drop
computed kinematically from the velocity triangles, $\Delta H_0 = c_p(T_{01}-T_{03})$,
must equal the value $\Delta H_{real}=P_{stage}/\dot m$ imposed by the power
balance. A Newton--Raphson loop adjusts $P_3$ until the two agree:
\[
\Delta H_0 < \Delta H_{real} \Rightarrow P_3' < P_3,\quad
\Delta H_0 > \Delta H_{real} \Rightarrow P_3' > P_3,\quad
\Delta H_0 = \Delta H_{real} \Rightarrow \text{OK}.
\]
The loop converges in 33 iterations to
$P_3 = \SI{1107265}{Pa}$, with $\Delta H_0 = \Delta H_{real} = \SI{219.27}{kJ/kg}$
(residual $< \SI{1e-5}{kJ/kg}$).

% =====================================================================
\section{Annulus Sizing}

Annulus sizing fixes the radial extent of the blade row. With a shrouded
end-wall assumption (equal hub and tip contour angles $\varepsilon$), the areas
follow from mass conservation, $A = \dot m/(\rho V_x)$, and the heights from
$h = A/(\pi D_m)$. The mean radius is set by the blade speed,
$R_m = 60\,U_m/(2\pi N) = \SI{1.11}{m}$. Table~\ref{tab:annulus} lists the
results.

\begin{table}[H]
\centering
\caption{Annulus geometry (constant mean radius).}
\label{tab:annulus}
\begin{tabular}{l r}
\toprule
\textbf{Quantity} & \textbf{Value}\\
\midrule
Mean radius $R_m$ [m]            & 1.110 \\
Span $h_2$ [mm]                  & 192.1 \\
Span $h_3$ [mm]                  & 234.6 \\
Tip radius (st.\ 3) $R_{T3}$ [m] & 1.228 \\
Hub radius (st.\ 3) $R_{H3}$ [m] & 0.994 \\
End-wall flare $\varepsilon$ [\degr] & 15.90 \\
$AN^2$ [m$^2$\,rpm$^2$]          & $1.341\times10^{7}$ \\
\bottomrule
\end{tabular}
\end{table}

The end-wall flare angle, obtained from the span growth over the axial chord,
$\varepsilon = \tan^{-1}\!\big((h_3-h_2)/c_x\big) = 15.90\degr$, stays below the
$15$--$20\degr$ range that bounds end-wall separation.

% =====================================================================
\section{Design Constraints}

Table~\ref{tab:constr} compares the obtained values against the aerodynamic and
mechanical design guidelines. Eight of the ten constraints are satisfied; two are
marginally violated -- the stator-exit Mach number $M_2 = 0.700$ sits exactly at
the lower edge of the $0.7$--$0.85$ band, and the span ratio $h_3/h_2 = 1.221$
slightly exceeds the recommended $1.2$. Both are borderline and could be cleared
by a small change of $\psi$, $\phi$ or $r$.

\begin{table}[H]
\centering
\caption{Comparison of design control parameters and obtained results.}
\label{tab:constr}
\begin{tabular}{l c c c}
\toprule
\textbf{Constraint} & \textbf{Design range} & \textbf{Obtained} & \textbf{Status}\\
\midrule
$M_2$              & $0.7$--$0.85$ & 0.700  & FAIL (low) \\
$\beta_2$ [\degr]  & $<47.5$       & 42.64  & PASS \\
$M_{w2}$           & $<0.5$        & 0.308  & PASS \\
$M_{w3}$           & $0.65$--$0.8$ & 0.666  & PASS \\
$\beta_3$ [\degr]  & $<75$         & 69.53  & PASS \\
$\Delta\beta$ [\degr] & $110$--$120$ & 112.17 & PASS \\
$M_3$              & $0.3$         & 0.300  & PASS \\
$\alpha_3$ [\degr] & $<35$         & 34.18  & PASS \\
$h_3/h_2$          & $<1.2$        & 1.221  & FAIL (high) \\
$AN^2$ [m$^2$rpm$^2$] & $<3\times10^7$ & $1.34\times10^7$ & PASS \\
\bottomrule
\end{tabular}
\end{table}
